\chapter{优化算法}

训练把经验风险写成可微标量 $f(\bm{x})$，再用迭代
$\bm{x}\leftarrow\bm{x}-\eta\,\widehat{\nabla f}(\bm{x})$ 更新参数。
深度学习要的是泛化，优化只直接压训练误差；二者目标不同。
下面先写凸性与梯度下降，再写随机、动量与自适应方法。

\section{优化和深度学习}

记训练目标为 $f(\bm{x})$。最大化等价于最小化 $-f$。
优化求 $f$ 的数值极小；深度学习还要求该极小对应的模型在未见数据上风险低。

\subsection{优化的目标}

经验风险与风险分别为
\[
R_{\mathrm{emp}}(\bm{x})
=
\frac{1}{n}\sum_{i=1}^{n}\ell\bigl(f_{\bm{x}}(\bm{z}^{(i)}),y^{(i)}\bigr),
\qquad
R(\bm{x})
=
\mathbb{E}_{(\bm{z},y)}[\ell(f_{\bm{x}}(\bm{z}),y)].
\]
优化算法直接减小 $R_{\mathrm{emp}}$；泛化要求 $R$ 小。
过拟合时 $R_{\mathrm{emp}}\ll R$，故仅压训练误差不够。

\subsection{深度学习中的优化挑战}

多数深度目标无解析解，只能用数值迭代。
梯度为零的点可以是局部极小、鞍点，或因梯度消失而停滞。

\subsubsection{局部最小值}

在 $[-1,2]$ 上
\[
f(x)=x\cos(\pi x)
\]
有多个局部极小。梯度下降停在哪一个，取决于初值与步长。

\subsubsection{鞍点}

$f(x)=x^3$ 在 $x=0$ 处 $f'(0)=f''(0)=0$，但不是极小。
$f(x,y)=x^2-y^2$ 在 $(0,0)$ 是鞍点：沿 $x$ 极小、沿 $y$ 极大。
Hessian 的特征值全正为局部极小，全负为局部极大，有正有负则为鞍点。

\subsubsection{梯度消失}

若 $f(x)=\tanh(x)$，则
\[
f'(x)=1-\tanh^2(x),\qquad f'(4)\approx 0.0013.
\]
从 $x=4$ 出发，步长 $\eta f'(x)$ 接近零，迭代长时间几乎不动。

\subsection{小结}

优化压 $R_{\mathrm{emp}}$，不自动压 $R$。
局部极小、鞍点与 $f'\approx 0$ 都会使
$x\leftarrow x-\eta f'(x)$ 停滞。

\subsection{练习}

核验：隐藏层宽度为 $d$ 时局部极小的置换对称、对称随机矩阵特征值分布是否关于零对称。

\section{凸性}

若算法在凸问题上失败，非凸上更难指望。
深度目标整体非凸，但局部常近似凸。

\subsection{定义}

先固定凸集，再固定凸函数。

\subsubsection{凸集}

集合 $\mathcal{X}$ 凸，当且仅当对任意 $a,b\in\mathcal{X}$ 与 $\lambda\in[0,1]$，
\[
\lambda a+(1-\lambda)b\in\mathcal{X}.
\]

\subsubsection{凸函数}

$f$ 在凸集上凸，当且仅当
\[
\lambda f(x)+(1-\lambda)f(x')
\geq
f\bigl(\lambda x+(1-\lambda)x'\bigr).
\]

\subsubsection{詹森不等式}

对 $\alpha_i\geq 0$、$\sum_i\alpha_i=1$，
\begin{align*}
\sum_i\alpha_i f(x_i)
&\geq
f\Bigl(\sum_i\alpha_i x_i\Bigr),
\\
\mathbb{E}_X[f(X)]
&\geq
f\bigl(\mathbb{E}_X[X]\bigr).
\end{align*}
取 $f=-\log$ 得
\[
\mathbb{E}_{Y\sim P(Y)}[-\log P(X\mid Y)]
\geq
-\log P(X).
\]

\subsection{性质}

凸性把局部信息变成全局信息。

\subsubsection{局部极小值是全局极小值}

设 $x^{\ast}$ 局部极小而 $f(x')<f(x^{\ast})$。对充分小的 $\lambda\in(0,1)$，
\begin{align*}
f\bigl(\lambda x^{\ast}+(1-\lambda)x'\bigr)
&\leq
\lambda f(x^{\ast})+(1-\lambda)f(x')
\\
&<
\lambda f(x^{\ast})+(1-\lambda)f(x^{\ast})
=f(x^{\ast}),
\end{align*}
与局部极小矛盾。故凸函数的局部极小即全局极小。

\subsubsection{凸函数的下水平集是凸的}

下水平集
\[
\mathcal{S}_b:=\{x\in\mathcal{X}:f(x)\leq b\}
\]
对凸 $f$ 是凸的：若 $x,x'\in\mathcal{S}_b$，则
\[
f\bigl(\lambda x+(1-\lambda)x'\bigr)
\leq
\lambda f(x)+(1-\lambda)f(x')
\leq b.
\]

\subsubsection{凸性和二阶导数}

一维时，凸性蕴含
\[
\frac12 f(x+\epsilon)+\frac12 f(x-\epsilon)
\geq f(x),
\]
从而
\[
f''(x)
=
\lim_{\epsilon\to 0}
\frac{f(x+\epsilon)+f(x-\epsilon)-2f(x)}{\epsilon^2}
\geq 0.
\]
若 $f'$ 单调，由中值定理
\[
f'(\alpha)=\frac{f(x)-f(a)}{x-a},\qquad
f'(\beta)=\frac{f(b)-f(x)}{b-x},
\]
再令 $\lambda=(x-a)/(b-a)$ 得
\[
\lambda f(b)+(1-\lambda)f(a)
\geq
f\bigl((1-\lambda)a+\lambda b\bigr).
\]
高维时定义
\[
g(z)\stackrel{\mathrm{def}}{=}
f\bigl(z\bm{x}+(1-z)\bm{y}\bigr),\qquad z\in[0,1].
\]
$f$ 凸当且仅当每个这样的 $g$ 凸，当且仅当 Hessian $\nabla^2 f$ 半正定。

\subsection{约束}

约束问题写成
\begin{align*}
\operatorname*{minimize}_{\bm{x}}&\quad f(\bm{x})
\\
\text{subject to}&\quad
c_i(\bm{x})\leq 0,\quad i=1,\ldots,N.
\end{align*}

\subsubsection{拉格朗日函数}

乘子 $\alpha_i\geq 0$，
\[
L(\bm{x},\alpha_1,\ldots,\alpha_n)
=
f(\bm{x})+\sum_{i=1}^{n}\alpha_i c_i(\bm{x}).
\]

\subsubsection{惩罚}

把 $\alpha_i c_i(\bm{x})$ 加进目标，近似强制 $c_i(\bm{x})\leq 0$。
权重衰减 $\frac{\lambda}{2}\|\bm{w}\|^2$ 对应约束 $\|\bm{w}\|^2-r^2\leq 0$。

\subsubsection{投影}

梯度裁剪
\[
\bm{g}\leftarrow \bm{g}\cdot\min\bigl(1,\theta/\|\bm{g}\|\bigr)
\]
把更新限制在半径 $\theta$ 的球上。一般投影为
\[
\operatorname{Proj}_{\mathcal{X}}(\bm{x})
=
\operatorname*{argmin}_{\bm{x}'\in\mathcal{X}}\|\bm{x}-\bm{x}'\|.
\]

\subsection{小结}

凸函数满足
$\lambda f(x)+(1-\lambda)f(x')\geq f(\lambda x+(1-\lambda)x')$，
局部极小即全局极小，Hessian 半正定。
约束经 $L$ 或惩罚进入目标，可行点由 $\operatorname{Proj}_{\mathcal{X}}$ 拉回。

\subsection{练习}

核验：$p$ 范数球的凸性、$\max(f,g)$ 保持凸而 $\min(f,g)$ 不必、Softmax 规范化的凸性。

\section{梯度下降}

全梯度下降很少直接用于深度模型，但它给出步长、曲率与预处理的原型。

\subsection{一维梯度下降}

泰勒展开
\[
f(x+\epsilon)=f(x)+\epsilon f'(x)+\mathcal{O}(\epsilon^2).
\]
取 $\epsilon=-\eta f'(x)$，
\[
f\bigl(x-\eta f'(x)\bigr)
=
f(x)-\eta f'(x)^2+\mathcal{O}\bigl(\eta^2 f'(x)^2\bigr).
\]
$\eta$ 充分小时 $f$ 下降，更新为
\[
x\leftarrow x-\eta f'(x).
\]

\subsubsection{学习率}

$\eta$ 过小则每步 $\lvert\eta f'(x)\rvert$ 太小；$|\eta f'(x)|$ 过大则二阶余项不可忽略，$f$ 不必下降。

\subsubsection{局部最小值}

$f(x)=x\cos(cx)$ 有无穷多个局部极小。过大的 $\eta$ 可把迭代送进较差的谷。

\subsection{多元梯度下降}

\[
\nabla f(\bm{x})
=
\biggl[
\frac{\partial f(\bm{x})}{\partial x_1},\ldots,
\frac{\partial f(\bm{x})}{\partial x_d}
\biggr]^{\mathsf T},
\]
\[
f(\bm{x}+\bm{\epsilon})
=
f(\bm{x})+\bm{\epsilon}^{\mathsf T}\nabla f(\bm{x})
+\mathcal{O}(\|\bm{\epsilon}\|^2),
\]
\[
\bm{x}\leftarrow \bm{x}-\eta\nabla f(\bm{x}).
\]

\subsection{自适应方法}

合适的 $\eta$ 难选。二阶方法用曲率自动定步，计算贵，但给出预处理的直觉。

\subsubsection{牛顿法}

\[
f(\bm{x}+\bm{\epsilon})
=
f(\bm{x})+\bm{\epsilon}^{\mathsf T}\nabla f(\bm{x})
+\frac12\bm{\epsilon}^{\mathsf T}\nabla^2 f(\bm{x})\bm{\epsilon}
+\mathcal{O}(\|\bm{\epsilon}\|^3).
\]
令 $\bm{H}=\nabla^2 f(\bm{x})$，驻点条件
\[
\nabla f(\bm{x})+\bm{H}\bm{\epsilon}=0
\qquad\Rightarrow\qquad
\bm{\epsilon}=-\bm{H}^{-1}\nabla f(\bm{x}).
\]

\subsubsection{收敛性分析}

一维牛顿误差 $e^{(k)}=x^{(k)}-x^{\ast}$ 满足
\begin{align*}
0&=f'\bigl(x^{(k)}-e^{(k)}\bigr)
\\
&=f'(x^{(k)})-e^{(k)}f''(x^{(k)})
+\frac12(e^{(k)})^2 f'''(\xi^{(k)}),
\end{align*}
从而
\[
\bigl\lvert e^{(k+1)}\bigr\rvert
=
\frac12(e^{(k)})^2
\frac{\lvert f'''(\xi^{(k)})\rvert}{f''(x^{(k)})}
\leq
c\,(e^{(k)})^2.
\]
在凸且 $f''$ 有下界时为二次收敛。

\subsubsection{预处理}

对角近似
\[
\bm{x}
\leftarrow
\bm{x}-\eta\,\operatorname{diag}(\bm{H})^{-1}\nabla f(\bm{x})
\]
按坐标缩放步长，避免直接求 $\bm{H}^{-1}$。

\subsubsection{梯度下降和线搜索}

沿 $-\nabla f(\bm{x})$ 对 $\eta$ 做一维搜索，使
$f(\bm{x}-\eta\nabla f(\bm{x}))$ 最小。
每步需在全数据上求值，深度模型上通常不可行。

\subsection{小结}

更新 $\bm{x}\leftarrow\bm{x}-\eta\nabla f(\bm{x})$ 中 $\eta$ 过大发散、过小停滞。
牛顿步 $-\bm{H}^{-1}\nabla f$ 在凸问题上快，非凸上 Hessian 不定时不能直接用。

\subsection{练习}

核验：不同 $\eta$ 下 $f$ 是否单调下降、线搜索是否需要导数、对角 Hessian 预处理对坐标尺度不均问题的作用。

\section{随机梯度下降}

全梯度每次用全部 $n$ 个样本；随机梯度每次用一个样本。

\subsection{随机梯度更新}

\[
f(\bm{x})=\frac1n\sum_{i=1}^{n}f_i(\bm{x}),
\qquad
\nabla f(\bm{x})=\frac1n\sum_{i=1}^{n}\nabla f_i(\bm{x}).
\]
随机更新
\[
\bm{x}\leftarrow \bm{x}-\eta\nabla f_i(\bm{x})
\]
无偏：
\[
\mathbb{E}_i[\nabla f_i(\bm{x})]=\nabla f(\bm{x}).
\]

\subsection{动态学习率}

\begin{align*}
\eta(t)&=\eta_i
&& t_i\leq t\leq t_{i+1}
&&\text{分段常数},
\\
\eta(t)&=\eta_0 e^{-\lambda t}
&&
&&\text{指数衰减},
\\
\eta(t)&=\eta_0(\beta t+1)^{-\alpha}
&&
&&\text{多项式衰减}.
\end{align*}

\subsection{凸目标的收敛性分析}

\[
\bm{x}_{t+1}
=
\bm{x}_t-\eta_t\partial_{\bm{x}}f(\bm{\xi}_t,\bm{x}),
\qquad
R(\bm{x})=\mathbb{E}_{\bm{\xi}}[f(\bm{\xi},\bm{x})].
\]
展开距离
\begin{align*}
\|\bm{x}_{t+1}-\bm{x}^{\ast}\|^2
&=
\|\bm{x}_t-\bm{x}^{\ast}\|^2
+\eta_t^2\|\partial_{\bm{x}}f(\bm{\xi}_t,\bm{x})\|^2
\\
&\quad
-2\eta_t\langle\bm{x}_t-\bm{x}^{\ast},\partial_{\bm{x}}f(\bm{\xi}_t,\bm{x})\rangle.
\end{align*}
若梯度有界 $\|\partial_{\bm{x}}f\|\leq L$，则
$\eta_t^2\|\partial_{\bm{x}}f\|^2\leq\eta_t^2 L^2$。
凸性给出
\[
f(\bm{\xi}_t,\bm{x}^{\ast})
\geq
f(\bm{\xi}_t,\bm{x}_t)
+\langle\bm{x}^{\ast}-\bm{x}_t,\partial_{\bm{x}}f(\bm{\xi}_t,\bm{x}_t)\rangle,
\]
于是
\begin{align*}
\|\bm{x}_t-\bm{x}^{\ast}\|^2
-\|\bm{x}_{t+1}-\bm{x}^{\ast}\|^2
&\geq
2\eta_t\bigl(f(\bm{\xi}_t,\bm{x}_t)-f(\bm{\xi}_t,\bm{x}^{\ast})\bigr)
-\eta_t^2 L^2.
\end{align*}
取期望并累加，加权平均
\[
\bar{\bm{x}}
\stackrel{\mathrm{def}}{=}
\frac{\sum_{t=1}^{T}\eta_t\bm{x}_t}{\sum_{t=1}^{T}\eta_t}
\]
满足
\[
\mathbb{E}[R(\bar{\bm{x}})]-R^{\ast}
\leq
\frac{r^2+L^2\sum_{t=1}^{T}\eta_t^2}{2\sum_{t=1}^{T}\eta_t}.
\]

\subsection{随机梯度和有限样本}

有放回时，一轮中某样本被抽到的概率
\[
1-(1-1/n)^n\approx 1-e^{-1}\approx 0.63.
\]
恰好抽一次的概率约为 $e^{-1}\approx 0.37$。
无放回遍历则每个样本每轮恰好一次。

\subsection{小结}

凸时 SGD 对广泛的 $\{\eta_t\}$ 收敛到最优；深度目标非凸，无同样保证。
$n$ 大时全梯度每步代价高，故用 $\nabla f_i$。

\subsection{练习}

核验：$\eta_t$ 衰减过快或过慢时 $\|\bm{x}_t-\bm{x}^{\ast}\|$ 的轨迹，以及有放回与无放回采样对无偏性的影响。

\section{小批量随机梯度下降}

全梯度数据效率低，单样本计算效率低。小批量在二者之间。

\subsection{向量化和缓存}

一次处理 $|\mathcal{B}|$ 个样本，使矩阵乘与缓存命中摊薄解释器与内存延迟。
多 GPU 时每卡至少一张图，有效批量随设备数放大。

\subsection{小批量}

单样本
\[
\bm{g}_t=\partial_{\bm{w}}f(\bm{x}_t,\bm{w}),
\]
小批量
\[
\bm{g}_t
=
\partial_{\bm{w}}
\frac{1}{|\mathcal{B}_t|}
\sum_{i\in\mathcal{B}_t}f(\bm{x}_i,\bm{w}).
\]
$\mathcal{B}_t$ 来自数据的随机排列，每轮每样本一次。

\subsection{读取数据集}

对已标准化的设计矩阵 $\bm{X}\in\mathbb{R}^{n\times d}$、响应 $\bm{y}$，
迭代器输出 $(\bm{X}_{\mathcal{B}},\bm{y}_{\mathcal{B}})$，
$|\mathcal{B}|$ 为批量。

\subsection{从零开始实现}

线性模型 $\hat{y}=\bm{w}^{\mathsf T}\bm{x}+b$，损失在批量上先平均再反传，
故更新中 $\bm{g}_t$ 不再除 $|\mathcal{B}|$。
$|\mathcal{B}|=n$ 退化为批量梯度下降。

\subsection{简洁实现}

框架优化器实现同一映射
$\bm{w}\leftarrow\bm{w}-\eta\bm{g}_t$，状态为空。

\subsection{小结}

小批量梯度
$\bm{g}_t=\partial_{\bm{w}}\frac{1}{|\mathcal{B}_t|}\sum_{i\in\mathcal{B}_t}f(\bm{x}_i,\bm{w})$
同时利用向量化与噪声梯度。训练中降低 $\eta$ 有助于后期收敛。

\subsection{练习}

核验：$|\mathcal{B}|$ 与 $\eta$ 如何改变目标下降速度，以及数据被复制一倍时三种梯度估计的偏差。

\section{动量法}

噪声梯度上 $\eta$ 难选。动量用过去梯度的泄漏平均代替当前梯度。

\subsection{基础}

先写平均，再写更新。

\subsubsection{泄漏平均值}

\[
\bm{g}_{t,t-1}
=
\partial_{\bm{w}}
\frac{1}{|\mathcal{B}_t|}
\sum_{i\in\mathcal{B}_t}f(\bm{x}_i,\bm{w}_{t-1}).
\]
\[
\bm{v}_t=\beta\bm{v}_{t-1}+\bm{g}_{t,t-1}
=
\sum_{\tau=0}^{t-1}\beta^{\tau}\bm{g}_{t-\tau,t-\tau-1}.
\]

\subsubsection{条件不佳的问题}

\[
f(\bm{x})=0.1\,x_1^2+2\,x_2^2
\]
两个方向曲率相差 $20$ 倍，固定 $\eta$ 的梯度下降在窄谷中振荡。

\subsubsection{动量法}

\begin{align*}
\bm{v}_t&\leftarrow\beta\bm{v}_{t-1}+\bm{g}_{t,t-1},
\\
\bm{x}_t&\leftarrow\bm{x}_{t-1}-\eta_t\bm{v}_t.
\end{align*}

\subsubsection{有效样本权重}

$\sum_{\tau=0}^{\infty}\beta^{\tau}=1/(1-\beta)$，
有效步长约为 $\eta/(1-\beta)$。
$\beta=0.9$ 相当于约 $10$ 个梯度的加权平均。

\subsection{实际实验}

速度 $\bm{v}$ 与 $\bm{g}$ 同形，需额外状态。

\subsubsection{从零开始实现}

维护 $\bm{v}$，按上式更新。增大 $\beta$ 时应略降 $\eta$，以免有效步长过大。

\subsubsection{简洁实现}

框架求解器在同一 $(\beta,\eta)$ 下实现同一 $\bm{v},\bm{x}$ 递推。

\subsection{理论分析}

二次凸目标把动量写成线性递推。

\subsubsection{二次凸函数}

\[
h(\bm{x})=\frac12\bm{x}^{\mathsf T}\bm{Q}\bm{x}+\bm{x}^{\mathsf T}\bm{c}+b.
\]
配平方后经 $\bm{Q}=\bm{U}\bm{\Lambda}\bm{U}^{\mathsf T}$ 变成
\[
h(\bm{z})=\frac12\bm{z}^{\mathsf T}\bm{\Lambda}\bm{z}+b'.
\]
无动量时
\[
\bm{z}_t=(\bm{I}-\bm{\Lambda})\bm{z}_{t-1}.
\]
有动量时
\begin{align*}
\bm{v}_t&=\beta\bm{v}_{t-1}+\bm{\Lambda}\bm{z}_{t-1},
\\
\bm{z}_t&=(\bm{I}-\eta\bm{\Lambda})\bm{z}_{t-1}-\eta\beta\bm{v}_{t-1}.
\end{align*}

\subsubsection{标量函数}

无动量 $x_{t+1}=(1-\eta\lambda)x_t$。有动量
\[
\begin{bmatrix}v_{t+1}\\ x_{t+1}\end{bmatrix}
=
\begin{bmatrix}
\beta & \lambda\\
-\eta\beta & 1-\eta\lambda
\end{bmatrix}
\begin{bmatrix}v_t\\ x_t\end{bmatrix}
=\bm{R}(\beta,\eta,\lambda)
\begin{bmatrix}v_t\\ x_t\end{bmatrix}.
\]
谱半径 $\rho(\bm{R})<1$ 时该坐标收敛。

\subsection{小结}

动量用 $\bm{v}_t=\sum_{\tau}\beta^{\tau}\bm{g}_{t-\tau}$ 替代当前梯度，
有效梯度个数为 $1/(1-\beta)$，并多存一个状态 $\bm{v}$。

\subsection{练习}

核验：$h(\bm{x})=\frac12\bm{x}^{\mathsf T}\bm{Q}\bm{x}+\bm{x}^{\mathsf T}\bm{c}+b$ 的最小点，以及 $\beta,\eta$ 对 $\rho(\bm{R})$ 的影响。

\section{AdaGrad算法}

稀疏特征上，全局递减的 $\eta$ 对常见坐标过慢、对罕见坐标过快。

\subsection{稀疏特征和学习率}

语言模型中罕见词的参数只在该词出现时更新。
若 $\eta_t=\mathcal{O}(t^{-1/2})$ 已很小，罕见坐标尚未被充分观测。

\subsection{预处理}

二次型经坐标缩放后
\[
f(\bm{x})=\bar f(\bar{\bm{x}})
=
\frac12\bar{\bm{x}}^{\mathsf T}\bm{\Lambda}\bar{\bm{x}}
+\bar{\bm{c}}^{\mathsf T}\bar{\bm{x}}+b,
\]
条件数 $\kappa=\Lambda_1/\Lambda_d$。对角预处理
\[
\tilde{\bm{Q}}
=
\operatorname{diag}^{-1/2}(\bm{Q})\,\bm{Q}\,\operatorname{diag}^{-1/2}(\bm{Q})
\]
把各坐标曲率拉近。梯度
\[
\partial_{\bar{\bm{x}}}\bar f(\bar{\bm{x}})
=
\bm{\Lambda}(\bar{\bm{x}}-\bar{\bm{x}}_0)
\]
幅度随坐标变化，可用 $|\bm{g}|$ 代理曲率。

\subsection{算法}

\begin{align*}
\bm{g}_t&=\partial_{\bm{w}}\ell\bigl(y_t,f(\bm{x}_t,\bm{w})\bigr),
\\
\bm{s}_t&=\bm{s}_{t-1}+\bm{g}_t^2,
\\
\bm{w}_t&=\bm{w}_{t-1}-\frac{\eta}{\sqrt{\bm{s}_t+\epsilon}}\cdot\bm{g}_t.
\end{align*}
平方与除法按坐标逐元素。测试函数仍用
$f(\bm{x})=0.1\,x_1^2+2\,x_2^2$。

\subsection{从零开始实现}

状态 $\bm{s}$ 与 $\bm{w}$ 同形，按上式更新。因 $\sqrt{\bm{s}_t}$ 单调增，后期步长自动变小，故可用比 SGD 更大的初始 $\eta$。

\subsection{简洁实现}

框架 AdaGrad 实现同一 $(\bm{s},\bm{w})$ 递推。

\subsection{小结}

AdaGrad 用 $\bm{s}_t=\sum_{\tau=1}^{t}\bm{g}_{\tau}^2$ 按坐标缩小步长：
大梯度坐标走得慢，稀疏坐标保持相对大的有效学习率。
深度非凸上 $\bm{s}_t$ 可能累积过快。

\subsection{练习}

核验：正交变换下 $\|\bm{U}\bm{c}-\bm{U}\bm{\delta}\|_2=\|\bm{c}-\bm{\delta}\|_2$，
以及旋转 $f(\bm{x})=0.1(x_1+x_2)^2+2(x_1-x_2)^2$ 后轨迹是否改变。

\section{RMSProp算法}

AdaGrad 的 $\bm{s}_t$ 无界累积，等价于 $\eta$ 按 $\mathcal{O}(t^{-1/2})$ 下降，非凸上往往过早停滞。

\subsection{算法}

\begin{align*}
\bm{s}_t&\leftarrow\gamma\bm{s}_{t-1}+(1-\gamma)\bm{g}_t^2,
\\
\bm{x}_t&\leftarrow\bm{x}_{t-1}-\frac{\eta}{\sqrt{\bm{s}_t+\epsilon}}\odot\bm{g}_t.
\end{align*}
展开
\[
\bm{s}_t
=(1-\gamma)\bigl(\bm{g}_t^2+\gamma\bm{g}_{t-1}^2+\gamma^2\bm{g}_{t-2}^2+\cdots\bigr).
\]
$\eta$ 由实验者单独调度；$\gamma$ 控制历史窗口，有效长度约 $1/(1-\gamma)$。

\subsection{从零开始实现}

同一二次函数上，AdaGrad 在 $\eta=0.4$ 后期几乎不动；RMSProp 因泄漏平均不让 $\bm{s}_t$ 无限增长。
深度网络常用 $\eta=0.01$、$\gamma=0.9$。

\subsection{简洁实现}

框架 RMSProp 实现同一更新。

\subsection{小结}

RMSProp 用泄漏平均 $\bm{s}_t\leftarrow\gamma\bm{s}_{t-1}+(1-\gamma)\bm{g}_t^2$ 做按坐标预处理，并把全局 $\eta$ 从累积平方和中分离。

\subsection{练习}

核验：$\gamma=1$ 时 $\bm{s}_t$ 不再更新的后果，以及旋转二次面上的收敛。

\section{Adadelta}

Adadelta 进一步取消显式 $\eta$，用参数变化量校准更新幅度。

\subsection{Adadelta算法}

梯度二阶矩
\[
\bm{s}_t=\rho\bm{s}_{t-1}+(1-\rho)\bm{g}_t^2.
\]
重缩放梯度与参数更新
\begin{align*}
\bm{g}_t'
&=
\frac{\sqrt{\Delta\bm{x}_{t-1}+\epsilon}}{\sqrt{\bm{s}_t+\epsilon}}\odot\bm{g}_t,
\\
\bm{x}_t&=\bm{x}_{t-1}-\bm{g}_t',
\\
\Delta\bm{x}_t
&=\rho\Delta\bm{x}_{t-1}+(1-\rho){\bm{g}_t'}^2.
\end{align*}
需两个状态 $\bm{s}_t$、$\Delta\bm{x}_t$。$\rho=0.9$ 对应约 $10$ 个半衰期。

\subsection{代码实现}

每坐标维护 $(\bm{s},\Delta\bm{x})$，按上式递推；框架 API 实现同一映射。

\subsection{小结}

Adadelta 用 $\sqrt{\Delta\bm{x}_{t-1}}/\sqrt{\bm{s}_t}$ 代替 $\eta$，仍用泄漏平均估计二阶统计。

\subsection{练习}

核验：能否不显式形成 $\bm{g}_t'$ 而得到同一 $\bm{x}_t$，以及与 AdaGrad、RMSProp 轨迹的差别。

\section{Adam算法}

Adam 把动量的一阶矩与 RMSProp 的二阶矩合在同一更新里，并做偏差校正。

\subsection{算法}

\begin{align*}
\bm{v}_t&\leftarrow\beta_1\bm{v}_{t-1}+(1-\beta_1)\bm{g}_t,
\\
\bm{s}_t&\leftarrow\beta_2\bm{s}_{t-1}+(1-\beta_2)\bm{g}_t^2.
\end{align*}
因 $\bm{v}_0=\bm{s}_0=\bm{0}$，
\[
\hat{\bm{v}}_t=\frac{\bm{v}_t}{1-\beta_1^t},\qquad
\hat{\bm{s}}_t=\frac{\bm{s}_t}{1-\beta_2^t}.
\]
\[
\bm{g}_t'=\frac{\eta\hat{\bm{v}}_t}{\sqrt{\hat{\bm{s}}_t}+\epsilon},
\qquad
\bm{x}_t\leftarrow\bm{x}_{t-1}-\bm{g}_t'.
\]

\subsection{实现}

状态为 $(\bm{v},\bm{s},t)$。常用 $\eta=0.01$，$\beta_1=0.9$，$\beta_2=0.999$。

\subsection{Yogi}

Adam 中 $\bm{s}_t$ 在梯度突变时可能被大正项拉动。把
\[
\bm{s}_t\leftarrow\bm{s}_{t-1}+(1-\beta_2)(\bm{g}_t^2-\bm{s}_{t-1})
\]
改成
\[
\bm{s}_t
\leftarrow
\bm{s}_{t-1}+(1-\beta_2)\bm{g}_t^2\odot\operatorname{sgn}(\bm{g}_t^2-\bm{s}_{t-1})
\]
使增量幅度受当前 $\bm{s}$ 限制。

\subsection{小结}

Adam 的更新是偏差校正后的
$\bm{x}_t\leftarrow\bm{x}_{t-1}-\eta\hat{\bm{v}}_t/(\sqrt{\hat{\bm{s}}_t}+\epsilon)$。
梯度尺度差异大时可用更大批量或 Yogi 的 $\bm{s}_t$。

\subsection{练习}

核验：去掉偏差校正后早期步长、以及何种 $\bm{g}_t$ 序列使 Adam 发散而 Yogi 收敛。

\section{学习率调度器}

更新规则确定方向，$\eta_t$ 确定步长。过大发散，过小停滞或次优；条件数大时更敏感。

\subsection{一个简单的问题}

在 Fashion-MNIST 上用固定 $\eta=0.3$ 训练修改后的 LeNet。
训练精度可继续升而测试精度停滞，间隙对应过拟合。

\subsection{学习率调度器}

调度器是映射 $t\mapsto\eta_t$。例如
\[
\eta_t=\eta_0(t+1)^{-1/2}.
\]
比固定 $\eta$ 更平滑，过拟合间隙往往更小。

\subsection{策略}

常用多项式、分段常数、余弦与预热。

\subsubsection{单因子调度器}

乘法衰减
\[
\eta_{t+1}\leftarrow\eta_t\cdot\alpha,\qquad\alpha\in(0,1),
\]
并加下限
\[
\eta_{t+1}\leftarrow\max(\eta_{\min},\eta_t\cdot\alpha).
\]

\subsubsection{多因子调度器}

在预定时刻 $s=\{5,10,20\}$ 上
$\eta_{t+1}\leftarrow\eta_t\cdot\alpha$。
直觉是：先用较大 $\eta$ 走到驻点附近，再减小 $\eta$ 降低参数噪声。

\subsubsection{余弦调度器}

\[
\eta_t
=
\eta_T+\frac{\eta_0-\eta_T}{2}\bigl(1+\cos(\pi t/T)\bigr).
\]

\subsubsection{预热}

先把 $\eta$ 从接近 $0$ 线性升到 $\eta_0$，再接入余弦或其他衰减，
避免深度网络初期过大步长导致发散。

\subsection{小结}

$\eta_t$ 应随训练下降。余弦调度为
$\eta_t=\eta_T+\frac{\eta_0-\eta_T}{2}(1+\cos(\pi t/T))$；
预热限制初期 $\|\eta\bm{g}\|$。

\subsection{练习}

核验：固定 $\eta$ 能达到的最优训练误差、多项式指数如何改变收敛，以及预热长度对初期损失的影响。
